\documentclass[11pt,a4paper]{article}

% --- Pakete ---
\usepackage[ngerman]{babel}
\usepackage[utf8]{inputenc}
\usepackage[T1]{fontenc}
\usepackage{lmodern} % Computer Modern (empfohlen)
\usepackage{geometry}
\geometry{
  a4paper,
  left=2.5cm,
  right=2.5cm,
  top=2.5cm,
  bottom=2.5cm
}
\usepackage{setspace}
\setstretch{1.24} % entspricht Empfehlung (1–1.5 Zeilenabstand)
\usepackage{microtype}
\usepackage{csquotes}
\usepackage{amsmath,amssymb}
\usepackage{graphicx}
\usepackage{caption}
\usepackage{subcaption}
\usepackage{siunitx}     % saubere Einheiten, Einheiten aufrecht
\sisetup{detect-all}
\usepackage{booktabs}
\usepackage{hyperref}
\usepackage{fancyhdr}

% --- Kopf- und Fußzeile ---
\pagestyle{fancy}
\fancyhf{}
\fancyhead[L]{Abschnitt \leftmark}
\fancyfoot[C]{\thepage}

% --- Nummerierte Gleichungen ---
\numberwithin{equation}{section}% --- Kopf- und Fußzeile ---
\pagestyle{fancy}
\fancyhf{}
\fancyhead[L]{Abschnitt \leftmark}
\fancyfoot[C]{\thepage}

% --- Nummerierte Gleichungen ---
\numberwithin{equation}{section}

% --- Abbildungs- und Tabellenbeschriftungen ---
\captionsetup{font=small,labelfont=bf}

% Kopfzeile: links aktueller Abschnitt
\fancyhead[L]{\nouppercase{\leftmark}}
\fancyfoot[C]{\thepage}

% Section-Mark definieren, damit \leftmark das Richtige enthält
\renewcommand{\sectionmark}[1]{%
  \markboth{\thesection\ #1}{}%
}


% --- Abbildungs- und Tabellenbeschriftungen ---
\captionsetup{font=small,labelfont=bf}

% --- Metadaten (für Titelseite) ---
\newcommand{\Lehrveranstaltung}{Praktikum Physikalische Eigenschaften der Materialien}
\newcommand{\Universitaet}{Universität Augsburg}
\newcommand{\Semester}{Wintersemester 2025/2026}
\newcommand{\Versuchsnummer}{1}
\newcommand{\Versuchstitel}{Brennstoffzelle}
\newcommand{\Gruppe}{C}
\newcommand{\Versuchsdatum}{02.03.2026}
\newcommand{\Abgabedatum}{TT.MM.JJJJ}

\newcommand{\Abgabeart}{Gemeinsames Versuchsprotokoll}
\newcommand{\Autor}{Fabio Nogielsky, Konstantin Szantho von Radnoth , Fabio Piskorz}

% ---------------------------------------------------
% Titelseite (ohne Seitenzahl)
% ---------------------------------------------------
\begin{document}
\pagenumbering{gobble} % keine Seitenzahl auf Titelseite

\begin{titlepage}
    \centering
    {\Large \Lehrveranstaltung\par}
    \vspace{0.5cm}
    {\large \Universitaet\par}
    \vspace{0.5cm}
    {\large \Semester\par}
    \vspace{1.5cm}
    {\LARGE \Abgabeart\par}
    \vspace{0.5cm}
    {\Large Versuchsnummer: \Versuchsnummer\par}
    \vspace{0.5cm}
    {\Large Versuchstitel: \Versuchstitel\par}
    \vspace{1.5cm}
    {\large Gruppe \Gruppe\par}
    \vspace{0.5cm}
    {\large Autor(en): \Autor\par}
    \vspace{1.5cm}
    {\large Versuchsdatum: \Versuchsdatum\par}
    \vspace{0.5cm}
    {\large Abgabedatum: \Abgabedatum\par}

    \vfill
\end{titlepage}

% ab hier Seitenzahlen normal
\pagenumbering{arabic}

% ---------------------------------------------------
% 1 Einleitung (max. ca. 1/2 Seite)
% ---------------------------------------------------
\section{Einleitung}
% Kurz und prägnant, max. 2/3 Seite. [file:1]
% - Warum ist der Versuch interessant/wichtig?
% - Welcher theoretische Hintergrund?
% - Was soll konkret untersucht werden?

Brennstoffzellen stellen einen vielversprechenden Ansatz für eine klimaneutrale Energieversorgung dar, da sie Wasserstoff elektrochemisch in elektrische Energie umwandeln und lediglich Wasser sowie Wärme als Reaktionsprodukte erzeugen. Im Rahmen der Energiewende gewinnen sie durch ihre hohe gravimetrische Energiedichte und Emissionsfreiheit an technischer Relevanz. Der Versuch behandelt das Funktionsprinzip von PEM-Elektrolyseur und PEM-Brennstoffzelle als zentrale Komponenten eines Wasserstoff-Energiesystems. Mittels U-I-Kennlinienaufnahmen und Wirkungsgradbestimmungen werden die Effizienzen der Energiewandler sowie irreversible Verluste durch Aktivierungs-, Ohmsche- und Konzentrationsüber-Spannungen experimentell charakterisiert. [file:1]
% ---------------------------------------------------
% 2 Theoretische Grundlagen (1–2 Seiten)
% ---------------------------------------------------
\section{Theoretische Grundlagen}
% Nur das, was für Verständnis und Auswertung nötig ist. [file:1]
% Herleitungen ggf. gekürzt und mit Literaturangaben.

In diesem Abschnitt werden die für den Versuch relevanten physikalischen Größen und Zusammenhänge erläutert. [file:1]

Als Beispiel für eine Gleichung mit erklärten Variablen:
\begin{equation}
    \Delta T = T_{\text{Raum}} - T_{\text{H2O}}.
\end{equation}
Hier bezeichnet \(T_{\text{Raum}}\) die Raumtemperatur und \(T_{\text{H2O}}\) die Wassertemperatur. [file:1]

% ---------------------------------------------------
% 3 Versuchsbeschreibung
% ---------------------------------------------------
\section{Versuchsbeschreibung}
% Versuchsaufbau und -durchführung so beschreiben,
% dass der Versuch ohne Anleitung nachvollziehbar ist. [file:1]

In diesem Abschnitt werden der Versuchsaufbau und die Versuchsdurchführung detailliert beschrieben, sodass der Versuch ohne zusätzliche Anleitung reproduzierbar ist. [file:1]

% Beispiel für eine Abbildung mit eindeutiger Nummer:
\begin{figure}[h]
    \centering
    %\includegraphics[width=0.7\textwidth]{aufbau.png}
    \caption{Schematische Darstellung des Versuchsaufbaus.}
    \label{fig:aufbau}
\end{figure}

Im Text wird auf die Abbildung \ref{fig:aufbau} verwiesen und deren Inhalt beschrieben. [file:1]

% ---------------------------------------------------
% 4 Auswertung und Diskussion
% ---------------------------------------------------
\section{Auswertung und Diskussion}
% Auswertung nach Aufgabenstellung, inkl. Fehlerrechnung
% und ggf. Vergleich mit Literaturwerten. [file:1]

In diesem Abschnitt werden die Messdaten ausgewertet, Unsicherheiten bestimmt und die Ergebnisse mit Literaturwerten verglichen. [file:1]

% Beispiel-Tabelle:
\begin{table}[h]
    \centering
    \caption{Beispielhafte Messwerte und berechnete Größen.}
    \label{tab:messwerte}
    \begin{tabular}{lccc}
        \toprule
        Messung & Zeit / \si{\second} & Weg / \si{\meter} & Geschwindigkeit / \si{\meter\per\second} \\
        \midrule
        1 & 1.0 & 0.50 & 0.50 \\
        2 & 2.0 & 1.00 & 0.50 \\
        \bottomrule
    \end{tabular}
\end{table}

Die Tabelle \ref{tab:messwerte} zeigt beispielhaft Messwerte mit zugehörigen Einheiten. [file:1]

% ---------------------------------------------------
% 5 Zusammenfassung (kurz, max. 1/2 Seite)
% ---------------------------------------------------
\section{Zusammenfassung}
% Kurzfassung des Versuchs und der wichtigsten Ergebnisse,
% typischerweise <= 1/2 Seite. [file:1]

Hier werden die wichtigsten Ergebnisse und Erkenntnisse des Versuchs kurz zusammengefasst und ggf. ein Ausblick gegeben. [file:1]

% ---------------------------------------------------
% 6 Laborbuch (Anhang)
% ---------------------------------------------------
\appendix
\section{Laborbuch}
% Digitale Kopie des Laborbuch-Eintrags (z.B. Scan oder Foto). [file:1]

Fügen Sie hier eine gut lesbare digitale Kopie der Laborbuchseiten ein, auf denen die Messdaten dokumentiert sind. [file:1]

% ---------------------------------------------------
% 7 Literaturverzeichnis
% ---------------------------------------------------
\section{Literaturverzeichnis}
% Einträge in der Reihenfolge des ersten Auftretens im Text. [file:1]
% Beispiel im IEEE-Stil:

\begin{thebibliography}{9}

\bibitem{Demtroeder1}
W.~Demtröder, \emph{Experimentalphysik 1: Mechanik und Wärme}, Springer, 2003. [file:1]

% weitere Quellen nach Bedarf

\end{thebibliography}

\end{document}
